\documentclass[a4paper]{article}
\usepackage[normalem]{ulem}
\usepackage{eurosym}
\usepackage[font=small,labelfont=bf]{caption}

% impostazioni generali
%Tutti gli usepackage vanno qui
\usepackage{geometry}
\usepackage[italian]{babel}
\usepackage[utf8]{inputenc}
\usepackage{tabularx}
\usepackage{longtable}
\usepackage{hyperref}
\usepackage{enumitem}
\usepackage{array} 
\usepackage{booktabs}
\newcolumntype{M}[1]{>{\centering\arraybackslash}m{#1}}
\usepackage[toc]{appendix}
\usepackage{caption}

\hypersetup{
	colorlinks=true,
	linkcolor=blue,
	filecolor=magenta,
	urlcolor=blue,
}
% Numerazione figure
\let\counterwithout\relax
\let\counterwithin\relax
\usepackage{chngcntr}

% distanziare elenco delle figure e delle tabelle
\usepackage{tocbasic}
\DeclareTOCStyleEntry[numwidth=3.5em]{tocline}{figure}% for figure entries
\DeclareTOCStyleEntry[numwidth=3.5em]{tocline}{table}% for table entries


%\counterwithout{table}{section}
%\counterwithout{figure}{section}
\captionsetup[table]{font=small,skip=5pt} 

\usepackage[bottom]{footmisc}
\usepackage{fancyhdr}
\setcounter{secnumdepth}{4}
\usepackage{amsmath, amssymb}
\usepackage{array}
\usepackage{graphicx}

\usepackage{ifthen}

\usepackage{float}
\restylefloat{table}

\usepackage{layouts}
\usepackage{url}
\usepackage{comment}
\usepackage{eurosym}

\usepackage{lastpage}
\usepackage{layouts}
\usepackage{eurosym}

\geometry{a4paper,top=3cm,bottom=4cm,left=2.5cm,right=2.5cm}

%Comandi di impaginazione uguale per tutti i documenti
\pagestyle{fancy}
\lhead{\includegraphics[scale=0.05]{../../../template/images/logo.png}}
%Titolo del documento
\rhead{\doctitle{}}
%\rfoot{\thepage}
\cfoot{Pagina \thepage\ di \pageref{LastPage}}
\setlength{\headheight}{41pt}
\setcounter{tocdepth}{5}
\setcounter{secnumdepth}{5}
\renewcommand{\footrulewidth}{0.4pt}

% multirow per tabelle
\usepackage{multirow}

% Permette tabelle su più pagine
%\usepackage{longtable}


% colore di sfondo per le celle
\usepackage[table]{xcolor}

%COMANDI TABELLE
\newcommand{\rowcolorhead}{\rowcolor[HTML]{007c95}}
\newcommand{\captionline}{\rowcolor[HTML]{FFFFFF}} %comando per le caption delle tabelle
\newcommand{\cellcolorhead}{\cellcolor[HTML]{007c95}}
\newcommand{\hlinetable}{\arrayrulecolor[HTML]{007c95}\hline}

%intestazione
% check for missing commands
\newcommand{\headertitle}[1]{\textbf{\color{white}#1}} %titolo colonna
\definecolor{pari}{HTML}{b1dae3}
\definecolor{dispari}{HTML}{d7f2f7}

% comandi \textit{Glossario}
\newcommand{\glo}{$_{G}$}
\newcommand{\glosp}{$_{G}$ }


%label custom
\makeatletter
\newcommand{\uclabel}[2]{%
	\protected@write \@auxout {}{\string \newlabel {#1}{{#2}{\thepage}{#2}{#1}{}} }%
	\hypertarget{#1}{#2}
}
\makeatother

%riportare pezzi di codice
\definecolor{codegray}{gray}{0.9}
\newcommand{\code}[1]{\colorbox{codegray}{\texttt{#1}}}

%versioni documenti
\newcommand{\AdRversione}{1.0.0}
\newcommand{\NdPversione}{1.0.0}
\newcommand{\PdPversione}{1.0.0}
\newcommand{\PdQversione}{1.0.0}
\newcommand{\Gloversione}{1.0.0}
\newcommand{\Stversione}{0.1.2}
\newcommand{\MUversione}{0.1.0}
%nome documento con versione 
\newcommand{\AdRdocumento}{\textit{Analisi dei Requisiti v.\AdRversione} }
\newcommand{\NdPdocumento}{\textit{Norme di Progetto v.\NdPversione} }
\newcommand{\PdPdocumento}{\textit{Piano di Progetto v.\PdPversione} }
\newcommand{\PdQdocumento}{\textit{Piano di Qualifica v.\PdQversione} }
\newcommand{\Glodocumento}{\textit{Glossario v.\Gloversione} }
\newcommand{\Stdocumento}{\textit{Glossario v.\Stversione} }
\newcommand{\MUDocumento}{\textit{Glossario v.\MUversione} }

% dati della prima pagina
% Configurazione della pagina iniziale
\newcommand{\doctitle}{\textit{Glossario}}
\newcommand{\rev}{0.3.1} % versione
\newcommand{\resp}{Ibra Elton} % inserire responsabile
\newcommand{\red}{Corbu Teodor Mihail}
\newcommand{\ver}{Andreetto Alessio}
\newcommand{\uso}{Esterno}
\newcommand{\dest}{\textit{Project Origin}
	\\ Prof. Vardanega Tullio 
	\\ Prof. Cardin Riccardo}
\newcommand{\describedoc}{Glossario contenente i termini per i quali è necessaria una definizione univoca del gruppo \newline \textit{Project Origin} nella realizzazione del progetto \textit{Personal Identity Wallet}}

 % per modificare la prima pagina editare questo file

\makeindex

\makeatletter
\renewcommand\paragraph{
\@startsection {paragraph}{4}{0mm}{-\baselineskip}{.5\baselineskip}{\normalfont \normalsize \bfseries }}
\makeatother

\begin{document}

% Prima pagina
\thispagestyle{empty}
\renewcommand{\arraystretch}{1.3}

\begin{titlepage}
	\begin{center}
		
	\includegraphics[scale = 0.20]{../../../template/images/logo.png}
	\\[1cm]
	\href{mailto:projectorigin2023@gmail.com}		      	
	{\large{\textit{projectorigin2023@gmail.com} } }\\[2.5cm]
	\Huge \textbf{\doctitle} \\[1cm]
	 \large
			 \begin{tabular}{r|l}
                        \textbf{Versione} & \rev{} \\
                        \textbf{Responsabile} & \resp{} \\
                        \textbf{Redattori} & \red{} \\ 
                        \textbf{Verificatori} &  \ver{} \\
                        \textbf{Uso} & \uso{} \\                        
                        \textbf{Destinatari} & \parbox[t]{5cm}{ \dest{} }
                \end{tabular} 
                \\[3.3cm]
                \large \textbf{Descrizione} \\ \describedoc{} 
     \end{center}
\end{titlepage}

% Diario delle modifiche
\section*{Registro delle modifiche}

\newcommand{\changelogTable}[1]{
	 

\renewcommand{\arraystretch}{1.5}
\rowcolors{2}{pari}{dispari}
\begin{longtable}{ 
		>{\centering}M{0.07\textwidth} 
		>{\centering}M{0.13\textwidth}
		>{\centering}M{0.20\textwidth}
		>{\centering}M{0.17\textwidth} 
		>{\centering\arraybackslash}M{0.30\textwidth} 
		 }
	\rowcolorhead
	\headertitle{Vers.} &
	\centering \headertitle{Data} &	
	\headertitle{Autore} &
	\headertitle{Ruolo} & 
	\headertitle{Descrizione} 
	\endfirsthead	
	\endhead
	
	#1

\end{longtable}
\vspace{-2em}

}

\changelogTable{
0.0.2 & 03-05-23 & Michele Beschin & Analista &  Aggiunti termini al glossario\\
0.0.1 & 02-05-23 & Michele Beschin & Analista &  Creazione struttura documento\\

} % editare questo
\pagebreak


% Indice
{
    \hypersetup{linkcolor=black}
    \tableofcontents
    \listoffigures %elenco figure
    \listoftables  %elenco tabelle
}
\pagebreak

% sezioni comuni 
% Informazioni generali
% \setcounter{table}{0} forse questo serve se l'indice delle tabelle parte da numeri sospetti
\section{Introduzione}

\subsection{Scopo de documento}
L’obiettivo del presente documento è di fornire una stima dei costi e dei tempi di consegna del progetto, nonché un piano dettagliato corredato di orari per ogni ruolo coinvolto. \\
La pianificazione sarà basata sull’analisi periodica delle attività svolte e una retrospettiva su quelle attività già completate, al fine di mantenere aggiornate e valide le 
stime effettuate. \\ Nel documento vengono inoltre identificate le possibili problematiche che il team potrebbe incontrare durante l’intero periodo di svolgimento del progetto.

\subsection{Scopo del prodotto}
Il capitolato prevede la realizzazione di un sistema di autenticazione dove un ente rilascia certificati di identità ad un utente, essi vengono memorizzati in un \textit{wallet}, 
per poi essere utilizzati per accedere a servizi. Sono previsti tre attori:
\begin{itemize}
    \item \textbf{Emittente}: entità che rilascia i certificati;
    \item \textbf{Holder}: persona fisica, che memorizza le credenziali di identità all'interno di un wallet;
    \item \textbf{Verifier}: entità che richiede delle credenziali per accedere a dei servizi.
\end{itemize}

\subsection{Glossario}
Al fine di evitare dubbi e ambiguità per determinati termini viene fornito un \textit{Glossario} che riporta la definizione specifica dei suddetti termini. 
Essi verranno contrassegnati con una lettera G (esempio\textsubscript{g}) a fine della parola. 
Tali termini verranno contrassegnati una sola volta per paragrafo onde evitare fastidiose ripetizioni. 

\subsection{Riferimenti}

\subsubsection{Riferimenti normativi}
\begin{itemize}
    \item \textit{Norme di progetto}: v 0.0.1
    \item Regolamento del progetto didattico: \\
    \url{https://www.math.unipd.it/~tullio/IS-1/2022/Dispense/PD02.pdf}
\end{itemize}

\subsubsection{Riferimenti informativi}
\begin{itemize}
    \item \textit{Analisi dei requisiti}: v 0.0.1
    \item Presentazione capitolato C3 \textit{Personal Identity Wallet} \\ \url{https://www.math.unipd.it/~tullio/IS-1/2022/Progetto/C3.pdf}
    \item Gestione di Progetto - slide T4 del corso Ingegneria del Software \\ \url{https://www.math.unipd.it/~tullio/IS-1/2022/Dispense/T04.pdf}
\end{itemize}


\newpage %forse si può togliere il newpage
\section{Qualità di Processo}
Per garantire la qualità dei processi il gruppo utilizza come riferimento lo standard \textbf{ISO/IEC 12207:1995}.
Dopo aver studiato i vari processi, si sono scelti quelli più adatti alle necessità del progetto.
\subsection{Processi primari}


\begin{longtable}{ 
    >{\centering}M{0.20\textwidth} 
    >{\centering}M{0.50\textwidth}
    >{\centering}M{0.17\textwidth} 
    }
\rowcolorhead
\headertitle{Processo} &
\centering \headertitle{Descrizione} &	
\headertitle{Metriche} 
\endfirsthead
\endhead

Fornitura & Scelta delle procedure e delle risorse necessarie per garantire che i servizi siano forniti in modo tempestivo & - \tabularnewline
Sviluppo & Realizzazione di un prodotto software di qualità, che soddisfi le esigenze del cliente& - \tabularnewline
\end{longtable}

\subsection{Processi di supporto}
\begin{longtable}{ 
    >{\centering}M{0.20\textwidth} 
    >{\centering}M{0.50\textwidth}
    >{\centering}M{0.17\textwidth} 
    }
\rowcolorhead
\headertitle{Processo} &
\centering \headertitle{Descrizione} &	
\headertitle{Metriche} 
\endfirsthead
\endhead

Verifica\glo & Si determina se i requisiti del prodotto sono soddisfatti& - \tabularnewline
Gestione della qualità & Viene garantita la conformità dei processi e la conformità con gli standard prefissati& - \tabularnewline
Documentazione & Controllo della leggibilità della documentazione prodotta& - \tabularnewline %per questo usiamo l'indice di gulpease
\end{longtable}
\subsection{Processi organizzativi}
\begin{longtable}{ 
    >{\centering}M{0.20\textwidth} 
    >{\centering}M{0.50\textwidth}
    >{\centering}M{0.17\textwidth} 
    }
\rowcolorhead
\headertitle{Processo} &
\centering \headertitle{Descrizione} &	
\headertitle{Metriche} 
\endfirsthead
\endhead

Gestione organizzativa & Controllo e organizzazione delle prestazioni di un processo\glo& - \tabularnewline
\end{longtable}
\section{Qualità di Prodotto}\newpage
\section{Resoconto attività di Verifica}
\subsection{Metriche di qualità di processo}
\begin{longtable}{ 
    >{\centering}M{0.20\textwidth} 
    >{\centering}M{0.30\textwidth}
    >{\centering}M{0.17\textwidth} 
    }
\rowcolorhead
\headertitle{Metrica} &
\centering \headertitle{Valore} &	
\headertitle{Esito} 
\endfirsthead
\endhead
%se non si riesce a testare qualche metrica, si indica "Non testato", e si lascia vuoto il campo valore
Planned Value & 12690 \euro& Superato \tabularnewline
Actual Cost & 12315 \euro& Superato \tabularnewline
Estimated At Completion & 12690 \euro& Superato \tabularnewline
Earned Value & 13289 \euro& Superato \tabularnewline
Estimated to Complete & 375 \euro& Superato \tabularnewline
Cost Variance & 974 \euro& Superato \tabularnewline
Schedule Variance & 3,15 \%& Superato \tabularnewline
Budget Variance & 375 \euro& Superato \tabularnewline %dovrebbe essere giusto
Requirements Stability Index & - & Non testato \tabularnewline
Code Coverage & - & Non testato \tabularnewline
Rischi non previsti & 0 & Superato \tabularnewline
Metriche Soddisfatte & 80\% & Superato \tabularnewline
\end{longtable}
\subsubsection{Indice di Gulpease}
\begin{longtable}{ 
    >{\centering}M{0.20\textwidth} 
    >{\centering}M{0.50\textwidth}
    >{\centering}M{0.17\textwidth} 
    }
\rowcolorhead
\headertitle{Documento} &
\centering \headertitle{Valore} &
\headertitle{Esito}
\endfirsthead
\endhead

Norme di Progetto & 70 & Superato \tabularnewline
Analisi dei Requisiti & 76 & Superato \tabularnewline
Piano di Progetto & 71 & Superato \tabularnewline
Piano di Qualifica & 69 & Superato \tabularnewline
Glossario & 72 & Superato \tabularnewline
Verbali Esterni & 70 & Superato \tabularnewline %per i verbali si è tenuto il valore più basso
Verbali Interni & 69 & Superato \tabularnewline
\end{longtable}

%\subsection{Metriche di qualità di prodotto}\newpage
\section{Test}
I test hanno verificato principalmente il corretto funzionamento del backend. Tramite il test degli endpoint è stato verificato il corretto flusso di funzionamento di tutte le componenti del backend, come il database.\\
Per il test del backend è anche stato utilizzato postman, per verificare manualmente, dove necessario, il test degli endpoint.\\
Il frontend è stato testato manualmente, verificando il corretto funzionamento delle pagine e dei componenti.\\

\subsection{Test di sistema}
%tabella (longtable) con 2 colonne: test e requisito
\begin{longtable}{ 
    >{\centering}M{0.10\textwidth} 
    >{\centering}M{0.70\textwidth} 
    >{\centering\arraybackslash}M{0.10\textwidth} 
    }
\rowcolorhead
\headertitle{\textbf{ID}} &
\headertitle{\textbf{Descrizione}} &
\centering \headertitle{\textbf{Esito}}
\endfirsthead
\endhead
TS01 & Verificato il corretto flusso di richiesta di una credenziale & PASSATO \tabularnewline
TS02 & Verificato il corretto flusso di approvazione di una richiesta di credenziale & PASSATO \tabularnewline
TS03 & Verificato il corretto flusso di rifiuto di una richiesta di credenziale & PASSATO \tabularnewline
TS04 & Verificato il corretto flusso di presentazione di una credenziale & PASSATO \tabularnewline
%aggiungerne altri
\end{longtable}
\subsection{Test di integrazione}%ce ne sono un paio, collegamenti tra i container
\begin{longtable}{ 
    >{\centering}M{0.10\textwidth} 
    >{\centering}M{0.70\textwidth} 
    >{\centering\arraybackslash}M{0.10\textwidth} 
    }
\rowcolorhead
\headertitle{\textbf{ID}} &
\headertitle{\textbf{Descrizione}} &
\centering \headertitle{\textbf{Esito}}
\endfirsthead
\endhead
TI01 & Collegamento tra Origin Issuer e Origin Wallet & PASSATO \tabularnewline
TI02 & Collegamento tra Origin Wallet e Origin Verifier & PASSATO \tabularnewline
\end{longtable}

\subsection{Test di unità}
\begin{longtable}{ 
    >{\centering}M{0.10\textwidth} 
    >{\centering}M{0.70\textwidth} 
    >{\centering\arraybackslash}M{0.10\textwidth} 
    }
\rowcolorhead
\headertitle{\textbf{ID}} &
\headertitle{\textbf{Descrizione}} &
\centering \headertitle{\textbf{Esito}}
\endfirsthead
\endhead
TU01 & Registrazione utente Origin Issuer & PASSATO \tabularnewline
TU02 & Login Utente Origin Issuer & PASSATO \tabularnewline
TU03 & Login Admin Origin Issuer & PASSATO \tabularnewline
TU04 & Richiesta Credenziale PID & PASSATO \tabularnewline
TU05 & Richiesta Credenziale Marital Status & PASSATO \tabularnewline
TU06 & Rilascio Credenziale con metodologia same device & PASSATO \tabularnewline
TU07 & Rilascio Credenziale con metodologia cross device & PASSATO \tabularnewline
TU08 & Richiesta Credenziale solo con token di autenticazione& PASSATO \tabularnewline
TU09 & Visualizzazione lista richieste di credenziale utente & PASSATO \tabularnewline
TU10 & Visualizzazione dettagli singola richista di credenziale PID & PASSATO \tabularnewline
TU11 & Visualizzazione dettagli singola richista di credenziale Marital Status & PASSATO \tabularnewline
TU12 & Visualizzazione lista richieste di credenziale admin & PASSATO \tabularnewline
TU13 & Approvazione richiesta di credenziale & PASSATO \tabularnewline
TU14 & Visualizzazione richieste di credenziale in pending & PASSATO \tabularnewline
TU15 & Visualizzazione richieste di credenziale not pending & PASSATO \tabularnewline
TU16 & Approvazione richieste di credenziale in pending & PASSATO \tabularnewline
TU17 & Rifiuto richieste di credenziale in pending & PASSATO \tabularnewline
TU18 & Registrazione utente Origin Wallet & PASSATO \tabularnewline
TU19 & Login utente Origin Wallet & PASSATO \tabularnewline
TU20 & Visualizzazione lista di credenziali utente & PASSATO \tabularnewline
TU21 & Visualizzazione dettaglio credenziale utente & PASSATO \tabularnewline
TU22 & Eliminazione credenziale utente& PASSATO \tabularnewline
TU23 & Presentazione credenziale & PASSATO \tabularnewline
TU24 & Approvazione presentazione credenziale& PASSATO \tabularnewline
TU25 & Rifiuto presentazione credenziale& PASSATO \tabularnewline
TU26 & Accesso ai servizi Verifier & PASSATO \tabularnewline

\end{longtable}
\pagebreak


\end{document}
